%% tags: [depth first search]
%% source: problems

\begin{prob}

    A modification of the \mintinline{python}{dfs} function from lecture is shown below.
    It is the same as the original \mintinline{python}{dfs} function, except that it has
    an added \mintinline{python}{print} statement.

    \begin{minted}[autogobble]{python}
        def dfs(graph, u, status=None):
            """Start a DFS at `u`."""
            # initialize status if it was not passed
            if status is None:
                status = {node: 'undiscovered' for node in graph.nodes}

            status[u] = 'pending'
            for v in graph.neighbors(u):
                if status[v] == 'undiscovered':
                    print("Hello!")
                    dfs(graph, v, status)
            status[u] = 'visited'
    \end{minted}

    Consider running this algorithm on the following graph, starting at node $a$:

    \includegraphics[width=.2\textwidth]{graph.png}

    How many times will \mintinline{python}{"Hello!"} be printed? Note that this is not a full DFS, and
    that not all nodes are reachable from $a$.

    \inlineresponsebox[.75in]{6}

\end{prob}
